% !TeX spellcheck = de_DE

Der eigentliche Verbesserungsprozess basiert auf einem iterativen, kontextbezogenen Re-Prompting-Ansatz.
Das verwendete Sprachmodell ändert sich im Laufe der Zeit, um die Konsistenz der Ergebnisse jedoch sicherstellen zu können, wird einheitlich für die Evaluation GPT-4 verwendet.

Werden im vorherigen Schritt bei der Bewertungslogik Verstöße gegen die Anforderungen festgestellt wird das dort dokumentierte Feedback in diesem Schritt weiterverwendet.
Das Feedback ist dabei beschreibend für die identifizierten Schwachstellen und dient als Kontext für die nächste Generierung.

Auf der Grundlage dieses Feedbacks erzeugt das Sprachmodell eine überarbeitete Version des Textes.
Anschließend wird die neue Version erneut bewertet.
Durch die wiederholte Anwendung dieses Verfahrens entsteht ein iterativer Optimierungsprozess, der auf die schrittweise Verbesserung der Textmetriken abzielt.

Die Überarbeitung orientiert sich nicht ausschließlich an den Fähigkeiten des Sprachmodells, sondern wird durch die Ergebnisse der automatisierten Qualitätsbewertung gesteuert.
Die verwendeten Metriken übernehmen somit die Rolle eines Kontrollmechanismus und einem Wegweiser für das LLM innerhalb des Feedbackloops.

Um eine unkontrollierte Anzahl von Iterationen zu vermeiden, wurden Abbruchkriterien definiert.
\begin{itemize}
	\item Der Verbesserungsprozess endet, sobald alle definierten Qualitätsanforderungen erfüllt sind.
	In diesem Fall wird davon ausgegangen, dass die angestrebte Textqualität erreicht wurde.
	\item Zusätzlich wird eine maximale Anzahl von Iterationen festgelegt.
	Wird diese Grenze erreicht ohne, dass sämtliche Anforderungen erfüllt werden, wird der aktuell beste Textstand als Ergebnis übernommen, jedoch mit einer Anmerkung, dass es zu einem Abbruch kam.
\end{itemize}